% !TeX spellcheck = nb_NO
% Kommentaren ovenfor gir instruks til editoren din om at den skal bruke norsk stavekontroll. Dette fungerer bl.a. i TeXstudio.

\documentclass[a4paper,norsk]{article}
% Norsk orddeling
\usepackage[norsk]{babel}
% AMS-pakkene er nyttige
\usepackage{amsmath,amsfonts,amssymb,amsthm}
% For å definere lister, som f.eks. deloppgaver
\usepackage{enumitem}
% Inneholder mye nyttig, slik som \coloneqq
\usepackage{mathtools}
% Penere kalligrafiske fonter
\usepackage[utf8]{inputenc}
\usepackage[T1]{fontenc}
\usepackage[a4paper, margin=0.5in]{geometry}
\usepackage{mathabx}
\usepackage{amsthm}
\usepackage{listings}
\usepackage{xcolor}
\usepackage{ifthen}
\usepackage{hyperref}

\usepackage[most]{tcolorbox}
\tcbuselibrary{theorems, breakable}
\tcbset{before upper=\par, after upper=\par} % allows itemize and enumerate inside boxes

\newtcbtheorem[]{definitionbox}{Definisjon}{
  colback=blue!3,
  colframe=blue!70!black,
  fonttitle=\bfseries\color{white},
  coltitle=white,
  boxrule=0.6pt,
  arc=2mm,
  enhanced,
  breakable
}{def}

\newtcbtheorem[]{theorembox}{Teorem}{
  colback=green!3,
  colframe=green!50!black,
  fonttitle=\bfseries,
  coltitle=black,
  boxrule=0.6pt,
  before upper=\par\ignorespaces,
  after upper=\par\unskip,
  arc=2mm,
  enhanced,
  breakable
}{thm}

\newtcbtheorem[]{oppgavebox}{Oppgave}{
  colback=gray!5,
  colframe=black!60,
  fonttitle=\bfseries,
  coltitle=black,
  boxrule=0.6pt,
  arc=2mm,
  enhanced,
  breakable
}{oppg}

\newtcolorbox{sectionbox}[1][]{
  colback=white,
  colframe=black!70,
  boxrule=0.7pt,
  arc=3mm,
  top=2mm, bottom=2mm, left=2mm, right=2mm,
  title=#1,
  enhanced,
  breakable
}

\newtcbtheorem[]{corollarybox}{Korollar}{
  colback=purple!5,
  colframe=purple!60!black,
  fonttitle=\bfseries\color{white},
  coltitle=white,
  boxrule=0.6pt,
  arc=2mm,
  enhanced,
  breakable,
  before upper=\par\ignorespaces,
  after upper=\par\unskip
}{cor}

\newtcbtheorem[]{lemmabox}{Lemma}{
  colback=orange!5,
  colframe=orange!70!black,
  fonttitle=\bfseries,
  boxrule=0.6pt,
  arc=2mm,
  enhanced,
  breakable
}{lem}

\newtcbtheorem[]{propositionbox}{Proposisjon}{
  colback=cyan!5,
  colframe=cyan!60!black,
  fonttitle=\bfseries,
  coltitle=black,
  boxrule=0.6pt,
  arc=2mm,
  enhanced,
  breakable
}{prop}

\newtcbtheorem[]{proofbox}{Bevis}{
  colback=white,
  colframe=black,
  fonttitle=\bfseries\color{white},
  coltitle=white,
  colbacktitle=black,
  title style={opacity=0.9},
  boxrule=0.6pt,
  arc=2mm,
  enhanced,
  breakable
}{proof}

\newlist{suboppgave}{enumerate}{1}
\setlist[suboppgave,1]{
  label=\textbf{(\alph*)},
  ref=\theenumi\alph*,
  leftmargin=1.5em,
  itemsep=0.4em,
  topsep=0.3em,
  parsep=0pt,
  partopsep=0pt
}



\setcounter{secnumdepth}{2}  % Sections and subsections numbered

% Python style  
\lstdefinestyle{pythonstyle}{
    language=Python,
    basicstyle=\ttfamily\small,
    keywordstyle=\color{blue},
    commentstyle=\color{green},
    numbers=left,
    frame=single
}

% Bruk bedre symboler for ≥, ≤, epsilon og phi
\renewcommand{\geq}{\geqslant}
\renewcommand{\leq}{\leqslant}
\newcommand{\eps} {\varepsilon}
\renewcommand{\epsilon}{\varepsilon}
\renewcommand{\phi}{\varphi}

% Definer de vanligste mengdene og funksjonene
\newcommand{\R}{\mathbb{R}}
\newcommand{\K}{\mathbb{K}}
\newcommand{\C}{\mathbb{C}}
\newcommand{\N}{\mathbb{N}}
\newcommand{\Z}{\mathbb{Z}}
\newcommand{\Q}{\mathbb{Q}}
% \L er definert som noe annet fra før av, så vi må redefinere \L
\renewcommand{\L}{\mathcal{L}}
% Rommet av polynomer
\newcommand{\Pol}{{\mathcal{P}}}
\DeclareMathOperator{\id}{id}
\DeclareMathOperator{\Image}{im}
\DeclareMathOperator{\Kernel}{ker}
\DeclareMathOperator{\Span}{span}
\DeclareMathOperator{\Rank}{rank}
\DeclareMathOperator{\Diag}{diag}
\DeclareMathOperator{\Proj}{proj}
\newcommand{\basis}{{\mathcal{B}}}
\newcommand{\basiss}{{\mathcal{C}}}
\newcommand{\basisss}{{\mathcal{D}}}

\usepackage[most]{tcolorbox}
\newtcolorbox{algoritme}[1]{
  title=Algoritme~#1,
  colback=white,
  colframe=black,
  boxrule=0.5pt,
  arc=2mm,
  left=3mm,
  right=3mm,
  top=1mm,
  bottom=1mm
}

% Definer oppgavemiljøet
\newenvironment{oppgave}[1]%
	{\trivlist{}\item\relax\textbf{Løsning på #1.}\medspace}%
	{\endtrivlist}
% Deloppgavelister
\newlist{deloppgave}{enumerate}{1}
\setlist[deloppgave,1]{label=(\textbf{\alph*}),align=left}

\title{Oblig 4 \\ MAT1125, høst 2025}
\author{Oliver Ekeberg}
\begin{document}
\maketitle

\tableofcontents

\section*{1.1 Vektorrom og underrom}
\addcontentsline{toc}{section}{1.1 Vektorrom og underrom}

\begin{definitionbox}{1.1.3}{}
La $ \mathbb{K}   $ være en kropp. Et \textbf{\textit{vektorrom over $ \mathbb{K} $ }}er en mengde U utstyrt med en operasjon + fra $ U \times U $ til U (vektoraddisjon) og en operasjon $ \cdot $ fra $ \mathbb{K} \times U $ til U (skalarmultiplikasjon) som tilfredsstiller

\begin{itemize}
    \item $ \left( u+v \right) + w = u + \left( v + w \right) \forall u, v, w \in U$ 
    \item $ u + v = v + u \forall u, v \in U $ 
    \item $ \alpha_{}^{} \cdot \left( u + v \right)  = \alpha_{}^{} u + \alpha_{}^{} v \forall \alpha_{}^{} \in \mathbb{K}$ og $ u, v \in U $
    \item $ \left( \alpha_{}^{} + \beta \right) u = \alpha_{}^{} u + \beta u \forall \alpha_{}^{} , \beta \in \mathbb{K} $ og $ u \in U $    
    \item $ \left( \alpha_{}^{} \beta \right) u = \alpha_{}^{} \left( \beta u \right)  \forall \alpha_{}^{} , \beta in \mathbb{K}  $ og $ u \in U $
    \item $ 1 \cdot u  = u \forall u \in U$
    \item $ \exists 0_U \in U $ s.a. $ u + \left( -u \right) =0_U  \forall u in U$     
\end{itemize}
\end{definitionbox}

\begin{propositionbox}{1.1.4}{}
La U være et vektorrom
\begin{itemize}
    \item $ \exists $ kun ett nullelement $ 0_U \in U $ 
    \item For enhver $ u \in U $ er elementet $ -u  $ unikt
    \item Om $ u + v = w $ er $ u = w - v $    
    \item $ \alpha_{}^{} \cdot 0_U = 0_U  \forall \alpha_{}^{} \in \mathbb{K}$ 
    \item $ 0 \cdot u = 0_U \forall u \in U$ 
    \item $ \left( -1 \right) \cdot u = -u \forall u \in  U $ 
\end{itemize}
\end{propositionbox}

\begin{proofbox}{1.1.4 i)}{}
Bevis på at det finnes kun ett nullelement

Anta først som motsigelse at $ u + v = u $ hvor  $ v \neq 0 $. Da er $ v = 0 $, men det motsiger antagelsen, så $ v = 0_U $  
\end{proofbox}

\begin{proofbox}{1.1.4 iv)}{}
Anta at $ \alpha_{}^{} \neq 0 $. Vi har at $ \alpha_{}^{} u = 0 $. Siden $ \alpha_{}^{} \neq 0 $ må $ u = 0 $. Men da er $ \alpha_{}^{} 0_U = 0_U $    
\end{proofbox}

\begin{definitionbox}{1.1.6 Underrom}{}
La U være et vektorrom over $ \mathbb{K} $. En ikke-tom delmengde av $ V \subseteq U $ er et \textbf{\textit{underrom}} av U dersom V er lukket under vektoraddisjon og skalarmultiplikasjon. dvs.
\begin{itemize}
    \item $ u + v \in V \forall u, v \in V $ 
    \item $ \alpha_{}^{} u \in V \forall u \in V  $ og $ \alpha_{}^{} \in  \mathbb{K} $  
\end{itemize}
\end{definitionbox}


\begin{oppgavebox}{1.1.7}{}
La 
\[
    U:= \mathbb{R}^2
\]
og $ V := \left\{ \begin{bmatrix}
    2r\\
    r
\end{bmatrix} : r \in \mathbb{R} \right\}  $ 

\begin{proofbox}{1.1.7}{}
V består av en vektor med en frihets grad. Det vil si alle vektorer i V er parallelle med $ \begin{bmatrix}
    2r\\
    r
\end{bmatrix}  $ 

la $ c \neq r $ og definer $ v_1 \in V  $ som $ v_1 = \begin{bmatrix}
    2c\\
    c
\end{bmatrix}  $   og $ v_2 \in V $ som $ v_2 = \begin{bmatrix}
    2r \\
    r
\end{bmatrix} $. Da er
\begin{align*}
    v_1 + v_2 &= \begin{bmatrix}
        2c\\
        c
    \end{bmatrix} + \begin{bmatrix}
        2r\\
        r
    \end{bmatrix} \\
    &= \begin{bmatrix}
        2 \left( c + r \right) \\
        \left( c+r \right) 
    \end{bmatrix} \in V
\end{align*}

Anta videre at vi har $ \alpha_{}^{} \in \mathbb{K} $ og definer $ \alpha_{}^{} v = \alpha_{}^{}  \begin{bmatrix}
    2r\\
    r
\end{bmatrix} = \begin{bmatrix}
    \alpha_{}^{} 2 r\\
    \alpha_{}^{} r
\end{bmatrix} \in V $  
\end{proofbox}
\end{oppgavebox}

\begin{oppgavebox}{1.1.8}{}
La U være et vektorrom
\begin{itemize}
    \item La V være et underrom av U. bevis at V i seg selv er et vektorrom
    \item La $ V_1, V_2 $ være underrom av U og la $ W := V_1 \cap V_2 $. Vis at W også er et underrom av U   
\end{itemize}

\begin{proofbox}{i)}{}
V i seg selv er et vektorrom fordi om $ u + v \in V  \forall u, v \in V$, så er også $ u + v \in U \forall u, v \in U $. Videre er også for en $ \alpha_{}^{} \in  \mathbb{K} $ om U er et vektorrom over $ \mathbb{K} $, så er $ \alpha_{}^{} u \in V  $ og $ \alpha_{}^{} u \in U \forall u \in  V$  
\end{proofbox}
\end{oppgavebox}


\section*{1.2 Lister av vektorer}
\addcontentsline{toc}{section}{1.2 Lister av vektorer}

\begin{definitionbox}{1.2.1}{}
La U være et vektorrom. En liste av vektorer er en tuppel $ \left( u_1, ..., u_n \right)  $ av vektorer $ u_1, ..., u_n \in U $. Om $ \mathcal{B}= \left( u_1, ..., u_n \right)  $ og $ \mathcal{C} = \left( v_1, ..., v_m \right)  $ er to lister av vektorer, sier vi at $ \mathcal{B} \subseteq \mathcal{C} $ om 
\[
    u_i \in \left\{ v_1, ..., v_m  \forall i = 1, ..., n \right\} 
\]
Da er $ \mathcal{C} $ en større liste enn $ \mathcal{B} $. Vi har ikke snakket noe om dimensjon ennå, men hver vektor i listen ligger i samme dimensjon som vektorrommet U
\end{definitionbox}

\section*{1.3 Lineær uavhengighet}
\addcontentsline{toc}{section}{1.3 Lineær uavhengighet}

Dette er et meget viktig kapittel. Lineær uavhegighet kan tenkes på som en måte å si at vi kan ikke uttrykke en vektor u uten vektoren u.\\

Ta frukt som et eksempel. En frukt kan ha forskjellige størrelser og former. Men vi kan ikke si at 
\[
    banan = eple + pære
\]
Uten å si at $ pære = \alpha_{}^{} \cdot banan $ eller $ eple = \alpha_{}^{} \cdot banan $ 
 
\begin{definitionbox}{1.3.1}{}
U vektorrom over $ K $. La $ \mathcal{B } = \left( u_1, ..., u_n \right)  $ være en liste av vektorer i U. En lineærkombinasjon av $ \mathcal{B} $ er et uttrykk på formen

\[
    \alpha_{1}^{} u_1 + ... + \alpha_{n}^{} u_n
\]
der $ \alpha_{1}^{}, ..., \alpha_{n}^{} \in \mathbb{K}  $. Spennet til $ \mathcal{B} $ er mengden av alle lineærkombinasjoner av $ \mathcal{B} $ 
\[
    span \mathcal{B} := \left\{ \alpha_{1}^{} u_1 + ... + \alpha_{n}^{} u_n : \alpha_{1}^{} , ..., \alpha_{n}^{} \in \mathbb{K} \right\} 
\]
Om $ V \subseteq U $ er en delmengde sier vi at $ \mathcal{B} $ utspenner V hvis $ Span \mathcal{B}  = V$ 
\end{definitionbox}

\begin{oppgavebox}{1.3.2}{}
U vektorrom over $ \mathbb{K} $ og $ \mathcal{B} = \left( u_1, ..., u_n \right)  $ en liste av vektorer i U. Vis at $ Span \mathcal{B}  $ er et underrom av U 

\begin{proofbox}{1.3.2}{}
Vi har
\[
    span \mathcal{B} := \left\{ \alpha_{1}^{} u_1 + ... + \alpha_{n}^{} u_n : \alpha_{1}^{} , ..., \alpha_{n}^{} \in \mathbb{K} \right\} 
\]
$ \forall u_j, u_k \in U $ så er $ \alpha_{j}^{} u_j + \alpha_{k}^{} u_k $ en lineærkombinasjon av vektorer i  $ span \left( \mathcal{B} \right)  $. Dermed er $ span \left( \mathcal{B} \right) \subseteq U  $ 

Videre om $ \alpha_{j}^{} \in \mathbb{K} $ , så er $ \alpha_{k}^{} u_k \in  span \left( \mathcal{B} \right)  $. Og for alle $ u_k \in span \left( \mathcal{B} \right)  $ , så er $ \alpha_{k}^{} u_k \in  U $ 
\end{proofbox}
\end{oppgavebox}

\begin{definitionbox}{1.3.3}{}
U vektorrom over $ \mathbb{K} $ og la $ \mathcal{B} = \left( u_1, ..., u_n \right)  $. Vi sier at $ \mathcal{B} $ er \textbf{\textit{lineært uavhengig}} hvis følgende implikasjon er sann $ \forall \alpha_{}^{} \in  \mathbb{K} $ 

\[
    \alpha_{1}^{} u_1 + ... + \alpha_{n}^{} u_n = 0 \rightarrow \alpha_{1}^{} = ... = \alpha_{n}^{} = 0
\]
$ \mathcal{B} $ er \textbf{\textit{lineært avhengig}} om $ \exists  \alpha_{1}^{} , ..., \alpha_{n}^{} \in \mathbb{K}$ som ikke alle er lik 0, og som er slik at 
\[
    \alpha_{1}^{} u_1 + ... + \alpha_{n}^{} u_n = 0
\]
Det dette sier er at i frukt eksemplet vårt, så er den eneste måte vi kan få ingen frukter i en mengde, er hvis 
\[
    \alpha_{1}^{}  banan + \alpha_{2}^{}  eple + \alpha_{3}^{}  pære = 0
\]
skjer hvis og bare hvis vi har $ \alpha_{1}^{}  = \alpha_{2}^{}  = \alpha_{3}^{} = 0  $. Vi kan godt si at $ 0 \cdot \left( frukt \right) = 0 $. Men nullelementet kan ikke være en vektor i en lineært uavhengig liste  
\end{definitionbox}

\begin{oppgavebox}{1.3.4}{}
La $ \mathbb{K} $ være en kropp. 
\begin{itemize}
    \item la $ e_1 := \begin{bmatrix}
        1\\
        0
    \end{bmatrix}  $  og $ e_2 := \begin{bmatrix}
        0\\
        1
    \end{bmatrix}  $ er lineært uavhengig i $ \mathbb{K}^2 $ 
    \item Generaliser forrige til listen $ \left( e_1, ..., e_n \right) \in \mathbb{K}^n  $ der $ e_j = \left( \delta_{jk} \right)_{k=1}^n  $.
\end{itemize} 

\begin{proofbox}{i)}{}
Vi setter at

\begin{align*}
    0 &= \alpha_{1}^{} e_1 + \alpha_{2}^{} e_2\\
    &= \alpha_{1}^{} \begin{bmatrix}
        1\\0
    \end{bmatrix} + \alpha_{2}^{} \begin{bmatrix}
        0\\1
    \end{bmatrix} \\
    &= \begin{bmatrix}
        \alpha_{1}^{} \\ \alpha_{2}^{} 
    \end{bmatrix} 
\end{align*}
Som har den eneste løsningen at $ \alpha_{1}^{} = \alpha_{2}^{} =0 $ 
\end{proofbox}

\begin{proofbox}{ii)}{}
Vi har at
\begin{align*}
    \alpha_{j}^{} e_j &= \left( \alpha_{j}^{} \delta_{jk} \right) 
\end{align*}
Av definisjon så er 

\[
    \left( \delta_{jk} \right) _{k=1}^n = \begin{cases}
        1, j=k\\
        0, j \neq k
    \end{cases}
\]
Derfor er $ 0 = \left( \alpha_{j}^{} \delta_{jk} \right)_{k=1}^n  $ sann hvis og bare hvis $ \alpha_{j}^{} =0 \forall j = 1, ..., n $ 
\end{proofbox}
\end{oppgavebox}

\begin{oppgavebox}{1.3.5}{}
La $ \mathcal{P}_1  = \left\{ \text{alle reelle polynomer av grad høyest lik 1} \right\} $ 
\begin{itemize}
    \item a) La $ p(t):= 1 , q(t) := t$. Vis at $ \left( p, q \right)  $ er lin. uavh.
    \item b) La $ p(t) = t+3 $, $ q(t) = t-3 $  
    \item c) La $ p(t) =t+1 $, $ q(t) = 3t+5 $ og $ r(t) = 2t - 4 $ . Vis at $ \left( p, q, r \right)  $ er linert AVHENGIG
\end{itemize}
\end{oppgavebox}


\begin{proofbox}{1.3.5 a)}{}
Vi har at $ p, q $ er lin. uavh. hvis og bare hvis 
\[
    0 = \alpha_{1}^{} p(t) + \alpha_{2}^{} q(t)
\]
er sann for $ \alpha_{1}^{} =\alpha_{2}^{} =0 $. 
\begin{align*}
    0 &= \alpha_{1}^{} p(t) + \alpha_{2}^{} q(t)\\
    &= \alpha_{1}^{} \cdot 1 + \alpha_{2}^{} \cdot t
\end{align*}
Dette skal gjelde for alle $ t \in  \mathbb{R} $. Vi ser at $ p(t) = 1 \neq 0  $ for noen $ t \in  \mathbb{R} $. Derfor kan uttrykket over bare være lik 0 hvis $ \alpha_{1}^{} = \alpha_{2}^{} =0 $ 
\end{proofbox}

\begin{proofbox}{1.3.5 b)}{}
Vi har
$$
  0 = \alpha_{1} p(t) + \alpha_{2} q(t) = \alpha_{1} \left( t+3 \right)  + \alpha_{2} \left( t-3 \right) 
$$

Her kan vi prøve forskjellige verdier for $t$. Anta først at $\alpha_{1}, \alpha_{2} \neq 0$.

For $t = 3$:
\[
  0 = \alpha_{1} \cdot 6 + \alpha_{2} \cdot 0
\]
er ikke sant. For $t = -3$:
\[
  0 = \alpha_{1} \cdot 0 - \alpha_{2} \cdot 6
\]
er heller ikke sant. Vi har altså at uttrykket er sant bare når $\alpha_{1} = \alpha_{2} = 0$.
\end{proofbox}

\begin{proofbox}{1.3.5 c)}{}
Vi har at 
\[
    0 = \alpha_{1}^{} \left( t + 1 \right) + \alpha_{2}^{} \left( 3t + 5 \right) + \alpha_{3}^{} \left( 2t -4  \right) 
\]

må være sann. Vi kan faktorisere ut og få etter noen steg at 

\begin{align*}
    0 &= t \left( \alpha_{1}^{} + 3 \alpha_{2}^{} + 2 \alpha_{3}^{}  \right) + \alpha_{1}^{} + 5 \alpha_{2}^{} - 4 \alpha_{3}^{} 
\end{align*}
For $ t=0 $, så får vi da at hele uttrykket er lik 0 om $ \alpha_{1}^{} + 5 \alpha_{2}^{} - 4 \alpha_{3}^{} =0 $  eller om $ \alpha_{1}^{} + 5 \alpha_{2}^{} = 4 \alpha_{3}^{}  $ 
Da er $ \alpha_{1}^{} p(t) + \alpha_{2}^{} q(t) + \alpha_{3}^{} r(t) = 0 $ for ikke alle $ \alpha_{1}^{}, \alpha_{2}^{}, \alpha_{3}^{} =0 $, og som fortsatt tilfredsstiller 
\[
    0 = \alpha_{1}^{} p(t) + \alpha_{2}^{} q(t) + \alpha_{3}^{} r(t)
\]
\end{proofbox}

\begin{oppgavebox}{1.3.6}{}
Vi sier at to vektorer er parallelle om $ u = \alpha_{}^{} v $. Forklar at dette er ekvivalent med å si at $ \left( u, v \right) $ er lineært avhengige
\end{oppgavebox}

\begin{proofbox}{1.3.6}{}
 Vi har at $ u = \alpha_{}^{} v $. Vi kan skrive om og få at 

\begin{align*}
    u &= \alpha_{}^{} v\\
    0 &= \alpha_{ }^{} v + \left( -1 \right) \cdot u \\
    0 &= \alpha_{1}^{} v + \alpha_{2}^{} u
\end{align*}
Dette er jo lik null, men $ \alpha_{2}^{} = -1 \neq 0 $ så det finnes skalarer $ \alpha_{1}^{}, \alpha_{2}^{}  $ som ikke alle er lik 0, men som fortsatt tilfredsstiller
\[
    0 = \alpha_{1}^{} v + \alpha_{2}^{} u
\]
Dermed er det ekvivalent å si at to vektorer er parallelle, og lineært avhengige
\end{proofbox}

\begin{oppgavebox}{1.3.7}{}
Oppgaven er
a) Vis at hvis $ \left( u_1, ..., u_n \right)  $ er lineært uavhengig, er ingen av vektorene $ u_1, ..., u_n $ lik null vektoren

b) Vis at følgende er ekvivalent med at $ \left( u_1, ..., u_n \right)  $ er lineært uavhengig: \\

Minst en av $ \alpha_{1}^{} , ..., \alpha_{n}^{}  $ er ulik 0 $ \rightarrow \alpha_{1}^{} u_1 + ... + \alpha_{n}^{} u_n \neq 0$   
c) Vis at følgende også er ekvivalent med at $ \left( u_1, ..., u_n \right)  $ er lineært avhengig:\\

En av vektorene  $ u_1, ..., u_n $ er en lineærkombinasjon av de andre
\end{oppgavebox}

\begin{proofbox}{1.3.7 a)}{}
Siden $ \left(  u_1, ..., u_n \right)  $ er lineært uavhengig, så er 
\[
    0 = \alpha_{1}^{} u_1 + ... + \alpha_{n}^{} u_n
\]
Sant hvis og bare hvis $ \alpha_{1}^{} = ... = \alpha_{n}^{} =0 $. Anta nå for motsigelse at $ 0_U $ er et element i listen $ u_1, ..., u_n $. Kall dette elementet for $ u_j $. Da vil tilhørende $ \alpha_{j}^{} = 1 $ gi at 
\[
    0 = \alpha_{1}^{} u_1 + ... + \alpha_{j}^{} u_j + ... + \alpha_{n}^{} u_n
\]
Hvor alle $ \alpha_{1}^{} , ..., \alpha_{j-1}^{} , \alpha_{j+1}^{} , ..., \alpha_{n}^{}   = 0$. Dette motstrider definisjonen for lineær uavhengighet, og nullementet kan derfor ikke være med i en lineært uavhengig liste.  
\end{proofbox}

\begin{proofbox}{1.3.7 b)}{}
Vi har en liste med lineært uavhengige vektorer $ u_1, ..., u_n $. Dette er et faktum vi har. 

Anta nå at vi har en mengde med koeffisienter $ \alpha_{1}^{}, ..., \alpha_{n}^{} $ hvor vi har minst en, si $ \alpha_{j}^{} \neq 0 $  for $ j \in \left\{ 1, ..., n \right\}  $. Da er dette ekvivalent med å si at 
\[
    \alpha_{1}^{} u_1 + ... + \alpha_{n}^{} u_n \neq 0
\]
Dette er fordi 
\begin{align*}
    \alpha_{1}^{} u_1 + ... + \alpha_{n}^{} u_n &= \alpha_{1}^{} u_1 + ... + \alpha_{j}^{} u_j + ... 0 \alpha_{n}^{} u_n \\
    &= \alpha_{j}^{} u_j
\end{align*}
som umulig kan være lik 0. Det er fordi vi har vist at nullelementet ikke kan være i en lineært uavhengig liste med vektorer, i tillegg til at vi har antatt at $ \alpha_{j}^{} \neq 0 $.\\

Anta motsatt at alle $ \alpha_{1}^{} , ..., \alpha_{n}^{} = 0 $. Da vil en  
\end{proofbox}

\begin{proofbox}{1.3.7 c)}{}
Vi har en liste med lineært avhengige vektorer $ \left( u_1, ..., u_n \right)  $. Da er det minst en $ \alpha_{j}^{} \neq 0 $ for $ j \in \left\{ 1, ..., n \right\}  $ slik at 

\[
    0 = \alpha_{1}^{} u_1 + ...+ \alpha_{j}^{} u_j + ...+ \alpha_{n}^{} u_n
\]
Dette er det samme som å skrive at 
\[
    - \alpha_{j}^{} u_j = \alpha_{1}^{} u_1 + ... + \alpha_{n}^{} u_n
\]
som sier at vektoren $ u_j $ er en lineær kombinasjon av de andre vektorene 
\end{proofbox}

\begin{lemmabox}{1.3.8}{}
La $ \left( u_1, ..., u_n \right)  $ være en liste av ikke-null vektorer i U. Da er følgende ekvivalent
\begin{itemize}
    \item $ \left( u_1, ..., u_n \right)  $ er lineært uavhengig
    \item $ u_k \notin span \left( u_1, ..., u_{k-1} \right)  $ for $ k=2, ..., n $ 
\end{itemize}
Hvis $ n=1 $ er begge utsagn automatisk sanne, derfor antar vi at $ n \geq 2 $ 
\end{lemmabox}

\begin{proofbox}{1.3.8}{}
Se i boken
\end{proofbox}

\begin{propositionbox}{1.3.9}{}
La $ \mathcal{B}= \left( u_1, ..., u_n \right)  $ være en liste an $  n \geq 2 $ vektorer i U, og anta at en av vektorene er en lin. komb. av de andre: 
\[
    u_k \in span \left( \mathcal{B}' \right) 
\]
der $ \mathcal{B}' := \left( u_1, ..., u_{k-1}, u_{k+1}, ..., u_n \right)  $. Da er

\[
    span \left( \mathcal{B}' \right) = span \left( \mathcal{B} \right) 
\]

Som et korollar gjelder det at for enhver liste med vektorer $ \mathcal{B} $  (uansett om den er lineært uanvhegig eller avhengig) slik at $ span \left( \mathcal{B} \right) \neq \left\{ 0 \right\}   $, finnes det en mindre liste $ \mathcal{C} \subseteq \mathcal{B} $ som er lineært uavhengig og som er slik at $ span \left( \mathcal{B} \right) = span \left( \mathcal{C} \right)   $ 

Se bevis i boken
\end{propositionbox}

\section*{1.4 Basis og dimensjon}
\addcontentsline{toc}{section}{1.4 Basis og dimensjon}

\begin{definitionbox}{1.4.1}{}
La U være et vektorrom. En \textbf{\textit{basis for U}} er en liste $ \left(  u_1, ..., u_n  \right)  $  av vektorer i U som er lineært uavhengig og som utspenner U
\end{definitionbox}


\begin{oppgavebox}{1.4.2}{}
Vis at standardbasisen $ \left( e_1, ..., e_n \right)  $ er en basis for $ \mathbb{K}^n $ 

\begin{proofbox}{1.4.2}{}
Vi har fra 1.3.4 at $ 0 = \alpha_{j}^{} e_j $ er sann hvis $ \alpha_{j}^{} = 0 \forall j \in \left\{ 1, ..., n \right\}  $. Derfor er $ \left( e_1, ..., e_n \right)  $ lineært uavhengig, og må bare vise at det utspenner $ \mathbb{K}^n $. 

Lar vi $ x = \left( x_j \right)_{j=1} \in \mathbb{K}^n  $ være vilkårlig, observerer vi at 
\[
    x = \sum_{ j=1 }^{ n } x_j e_j
\]
Altså ligger x i hele spennet til listen. Og siden x var vilkårlig, spenner hele listen $ \mathbb{K}^n $ 
\end{proofbox}
\end{oppgavebox}

\begin{definitionbox}{1.4.3}{}
Et vektorrom er \textbf{\textit{endeligdimensjonalt}} hvis det utspennes av endelig mange vektorer., ellers er det \textbf{\textit{uendeligdimensjonalt}}. Det trivielle vektorrommet $ \left\{ 0 \right\}  $ defineres som endeligdimensjonalt

Sagt med andre ord, Et vektorrom U er endeligdimensjoalt om $ \exists $ en liste $ \left( u_1, ..., u_n \right)  $ av vektorer i U slik at $ U = span \left( u_1, ..., u_n \right)  $ 
\end{definitionbox}

\begin{oppgavebox}{1.4.4}{}
La U være et ikke-trivielt vektorrom. Bevis at følgende er ekvivalent
\begin{itemize}
    \item U er endeligdimensjonalt
    \item U har en basis
\end{itemize}

\begin{proofbox}{1.4.4}{}
Anta først at U er endeligdimensjonalt. Det vil si at $ \exists  $ en liste $ \left( u_1, ..., u_n \right)  $ av vektorer i U slik at $ U = span \left( u_1, ..., u_n \right)  $ dvs. at en vektor $ u \in U $ kan skrive som
\[
    u = \alpha_{1}^{} u_1 + ... + \alpha_{n}^{} u_n
\]
Dette er en måte å si at basisen til U er listen $ \left( u_1, ..., u_n \right)  $ 

Skal nå vise at hvis U har en basis, så er U endeligdimensjonalt

Hvis U har en basis $ \left( u_1, ..., u_n \right)  $ så er listen $ \left( u_1, ..., u_n \right) $ lineært uavhengig, og dermed er en vilkårlig vektor $ u \in U $ uttrykt som
\[
    u = \alpha_{1}^{} u_1 + ... + \alpha_{n}^{} u_n
\]
Siden dette gjelder for en vilkårlig $ u \in U $ så er hele spennet til basisen lik U. Spennet til U er alle mulige lineærkombinasjoner av basisvektorene til et vektorrom U
\end{proofbox}
\end{oppgavebox}

\begin{theorembox}{1.4.5}{}
La U være et vektorrom over $ \mathbb{K} $ med en basis $ \mathcal{B} = \left( u_1, ..., u_n \right)  $. For enhver $ u \in U $ finnes da en unik n-tupper $ \left( x_1, ..., x_n \right) \in \mathbb{K}^n  $ slik at 
\[
    u = x_1 u_1 + ... + x_n u_n
\]
Vi betegner da denne n-tuppelen $ \left[ u \right] _{\mathcal{B}} = \left( x_1, ..., x_n \right)  $ som kalles \textbf{\textit{basisrepresentasjonen av u i basisen $ \mathcal{B} $ }}
\end{theorembox}

\begin{proofbox}{1.4.5}{}
Siden $ \mathcal{B} $ er en basis, så er $ \left( u_1, ..., u_n \right)  $  lineært uahengig. Anta at vi nå har en n-tuppel $ \left( x_1, ..., x_n \right)  $ som er slik at $ u = x_1 u_1 + ... + x_n u_n $. Anta vi har også en annen n-tuppel $ \left( x_1', ..., x_n' \right)  $  som er slik at $ u = x_1' u_1 + ... + x_n' u_n $. Da vil

\begin{align*}
    u - u &= \left( x_1 - x_1' \right) u_1 + ... + \left( x_n - x_n' \right) u_n 
\end{align*}
Siden $ \left( u_1, ..., u_n \right)  $ er lineært uavhengig, så må $ x_1=x_1', ..., x_n = x_n' $ og representasjonen $ u = x_1 u_1 + ... + x_n u_n $ er unik
\end{proofbox}

\begin{oppgavebox}{1.4.6}{}
a) La $ \mathcal{C} = \left( e_1, ..., e_n \right) $ være standardbasisen i $ \mathbb{K}^n $. Vis at
\[
    \left[ x \right] _{\mathcal{C}} = x \forall x \in \mathbb{K}^n
\]\\


b) La U være et vektorrom over $ \mathbb{K} $ med en basis $ \mathcal{B} = \left( u_1, ..., u_n \right)  $. Vis at 
\[
    \left[ u_j \right] _{\mathcal{B} } = e_j \forall j=1, ..., n
\]    
\end{oppgavebox}


\begin{proofbox}{1.4.6 a)}{}
For enhver $ x \in \mathbb{K}^n $ har vi at x er på formen


\[
    x = \begin{bmatrix}
        x_1 \\
        \vdots\\
        x_n
    \end{bmatrix} 
    = \begin{bmatrix}
        x_1 \\
        \vdots\\
        0
    \end{bmatrix} e_1 + ... + \begin{bmatrix}
        0 \\
        \vdots\\
        x_n
    \end{bmatrix} e_n = x_1 e_1 + ... x_n e_n 
\]
dermed er $ \left[ x \right] _{\mathcal{C}} = \left( x_1, ..., x_n \right) $ 
\end{proofbox}

\begin{proofbox}{1.4.6 b)}{}
Vi har at for en vilkårlig $ u_j \in U $, så er 
\[
    u_j = 0 \cdot u_1 + ... + 1 \cdot u_j + ... + 0 \cdot u_n
\] 

Da er 
\[
    \left[ u_j \right] _{\mathcal{B}} = \left( 0, ..., 1_j, ..., 0 \right) = e_j
\]
 
\end{proofbox}

\begin{lemmabox}{1.4.7}{}
La $ \left( u_1, ..., u_n \right)  $ og $ \left( v_1, ..., v_m \right)  $ være lister av vektorer i U. Anta at $ u_1, ..., u_n $ utspenner U, og $ \left( v_1, ..., v_m \right)  $ er lineært avhengig- Da er $ m \leq n $ 
\end{lemmabox}
\begin{proofbox}{1.4.7}{}
Se boken for bevis
\end{proofbox}

\begin{theorembox}{1.4.8}{}
La U være et endeligdimensjonalt vektorrom. Da er alle basiser for U like store. dvs. de inneholder like mange vektorer
\end{theorembox}

\begin{proofbox}{1.4.8}{}
Anta at U har to basiser $ \mathcal{C} = \left( v_1, ..., v_m \right)  $ og $ \mathcal{B} = \left( u_1, ..., u_n \right)  $. Skal vise at $ m \leq n $ og at $ m\geq n $. Anta først at $ \left( u_1, ..., u_n \right) $ utspenner U, og at $ \mathcal{C} $ er en lineært uavhengig liste i U. Da er ifølge lemmabox 1.4.7 $ m \leq n $. Anta nå motsatt at $ \left( v_1, ..., v_m  \right)  $ utspenner U og at $ \left( u_1, ..., u_n \right)  $ er lineært uavhengig. Da er $ m \leq n $. Og jeg har da vist at det eneste alternativet, er at $ m=n $  
\end{proofbox}

\begin{definitionbox}{1.4.9}{}
La U være et vektorrom over $ \mathbb{K}$. His U er endeligdimensjonal, sier vi at \textbf{\textit{dimensjonen til U }} er antall vektorer i U. Når $ U = \left\{ 0 \right\}  $ definerer vi $ dim U := 0 $. Dimensjonen til et uendeligdimensjonalt vektorrom Y definerer vi som $ dim U := \infty $ 
\end{definitionbox}

\begin{oppgavebox}{1.4.10}{}
a) La $ \mathbb{K} $  være en kropp og la $ n \in \mathbb{N} $. Vis at dimensjonen til $ \mathbb{K}^n $ over $ \mathbb{K} $ er n

b) Som vi så i 1.1.5 d), kan ethvert vektorrom over $ \mathbb{C} $ også ses på som et vektorrom over $ \mathbb{R} $. Vis at dimensjonen til $ \mathbb{C}^n $ er 2n


\begin{proofbox}{1.4.10 a)}{}
Vi har at en basis for $ \mathbb{K}^n$  er $ \left( e_1, ..., e_n \right)  $. Siden denne basisen inneholder n vektorer er
\[
    dim \mathbb{K}^n = n
\]
\end{proofbox}

\begin{proofbox}{1.4.10 b)}{}
Et tall i $ \mathbb{C} $ kan skrives på formen $ a + ib $ hvor $ a, b \in \mathbb{R} $. Det vil si at enhver vektor i $ \mathbb{C} $ må besrkives med to vektorer. Siden $ dim \mathbb{R}^n = n $ er dimensjonen til $ \mathbb{R}^n $ (se forrige oppgave, hvor $ dim \mathbb{K}^n = n $ 
\end{proofbox}
\end{oppgavebox}

\begin{propositionbox}{1.4.11}{}
La U være et vektorrom og V et underrom av U. Da er $ dim V \leq dimU $. Om U er endeligdimensjonal er $ dim V = dim U $ hvis og bare hvis $ V=U $ 
\end{propositionbox}

\begin{corollarybox}{1.4.12}{}
La $ n \in \mathbb{N} $, la U være et n-dim vektorrom, og la $ \mathcal{B} := \left( u_1, ..., u_n \right)  $ være en liste av n vektorer i U. Da er følgende ekvivalent
\begin{itemize}
    \item $ \mathcal{B}  $ er en basis for U
    \item $ \mathcal{B} $ er lineært uavhengig
\end{itemize} 
\end{corollarybox}


\section*{1.5 Polynombasiser}
\addcontentsline{toc}{section}{1.5 Polynombasiser}

Vi har at mengden av polynomer med grad høyest lik n er 
\[
    \mathcal{P}_n = \left\{ \text{alle reelle polynomer av grad høyest lik n} \right\} 
\]

\subsection*{1.5.1 Den kanoniske basisen}
\addcontentsline{toc}{section}{Den kanoniske basisen}

\begin{propositionbox}{1.5.1}{}
Definer $ p_k : \mathbb{R} \rightarrow \mathbb{R} $ som
\begin{itemize}
    \item $ p_0(t) = 1 $
    \item $ p_k(t) = t^k $  
\end{itemize}
$ \forall t \in  \mathbb{R}, k =1, ..., n $
Da er $ \left( p_0, ..., p_n \right)  $  en basis for $ \mathcal{P}_n $ 
\end{propositionbox}

\begin{oppgavebox}{1.5.2}{}
Vis at 
\begin{itemize}
    \item $ dim \mathcal{P}_n = n+1 $
    \item $ dim \mathcal{P} = \infty $  
    \item $ dim C^0 \left( \mathbb{R}, \mathbb{R} \right) = \infty $ 
\end{itemize}

\begin{proofbox}{1.5.2 a)}{}
Vektorrommet $ \mathcal{P}_n $ kan representeres med basisen $ \left( p_0, p_1, ..., p_n \right)  $ som har $ n+1 $ elementer i seg. Ergo er $ dim \mathcal{P}_n = n+1$ 
\end{proofbox}

\begin{proofbox}{1.5.2 b)}{}
Vektorrommet $ \mathcal{P} = \left\{ \text{alle reelle polynomer} \right\}  $. Vi kan begynne med å vise at $ p_{k+1} \notin span \left( p_0, ..., p_k \right)  $. La $ \left( p_0, ..., p_n \right)  $ være den kanoniske basisen, da er $ p_n $ et polynom av grad n, og $ p_{n+1} $ er et polynom av grad n+1. Dette kan ikke uttrykkes som en kombinasjon av polynomer av høyst grad n
\[
    p_{n+1} \notin span \left( p_0, ..., p_n \right) 
\]
Dette gjelder $ \forall n \geq 2 $ 
\end{proofbox}


\begin{proofbox}{1.5.2 c)}{}
$ C^0 \left( \mathbb{R}, \mathbb{R} \right)  $ er rommet av alle kontinuerlige funksjoner som er på formen $ f: \mathbb{R} \rightarrow \mathbb{R} $. Vi vet at $ \mathcal{P} $ er et underrom av dette rommet, fordi alle polynomer er kontinuerlige funksjoner. bevis for dette: Velg en $ u \in C^0 \left( \mathbb{R}, \mathbb{R} \right)  $ som er slik at $ u \notin \mathcal{P} $ og det er klart at $ \mathcal{P} $ er et underrom av $ C^0 \left( \mathbb{R}, \mathbb{R} \right)  $, fordi vektoraddisjon eller skalarmultiplikasjon av to vektorer i $ \mathcal{P} $ aldri kan forlate rommet $ \mathcal{P} $ 
\end{proofbox}
\end{oppgavebox}

\subsection*{1.5.2 Newton-basisen}
\addcontentsline{toc}{section}{Newton-basisen}

\begin{propositionbox}{1.5.3}{}
La $ t_0, ..., t_{n-1} \in \mathbb{R} $ være n forskjellige tall, og definer $ p_k : \mathbb{R} \rightarrow \mathbb{R} $ som
\[
    p_0(t) = 1
\]
\[
    p_k(t) := \prod_{j=0}^{k-1} \left( t - t_j \right)  \forall t \in R, k = 1, ..., n
\]
Da er $ \left( p_0, ..., p_n \right)  $ en basis for $ \mathcal{P}_n $ 
\end{propositionbox}

Polynomet av grad n som interpolerer en funksjon i $ n+1 $ forskjellige punkter $ t_0, ..., t_n $ er 
\[
    p(t) := \alpha_{0}^{} + \alpha_{1}^{} \left( t - t_1 \right) + \alpha_{2}^{} \left( t - t_1 \right) \left( t - t_2 \right) + ... + \alpha_{n}^{} \left( t- t_1 \right) ... \left( t - t_n \right) 
\]

der $ \alpha_{k}^{} = f \left[ t_0, ..., t_n \right]  $ 
der 
\[
    \begin{cases}
        f \left[ t_i \right] := f(t_i)\\
        f \left[ t_i, ..., t_j \right] := \frac{ f \left[ t_{i+1}, ..., t_j \right] - f \left[ t_i, ..., t_{j-1} \right] }{ t_j - t_i }
    \end{cases}
\]
Som et eksempel for $ n=2 $ får vi

\[
    p(t) = f(t_0) + \left( t-t_0 \right) \frac{ f(t_1) - f(t_0) }{ t_1 - t_0 } + \left( t - t_0 \right) \left( t - t_1 \right) \frac{ \frac{ f(t_2) - f(t_1) }{ t_2 - t_1 } - \frac{ f(t_1) - f(t_0) }{ t_1 - t_0 } }{ t_2 - t_0 }
\]
Det viktigste å trekke fra dette er at et polynom på formen $ p_n $ er et polynom av grad n. dvs. at det høyeste eksponenten vi finner er $ t^n $ 

\subsection*{1.5.3 Lagrange polynomer}
\addcontentsline{toc}{section}{Lagrange polynomer}

La $ t_0, .., t_{n} $ være $ n+1 $ forskjellige reelle tall. De tilhørende lagrange polynomene er funksjonene
\[
    L_k(t) := \frac{ \prod_{j\neq k} \left( t - t_i \right)  }{ \prod_{j\neq k}  \left( t_k - t_j  \right) } \forall t \in \mathbb{R}, k=0, ..., n
\]
Merk at 
\begin{itemize}
    \item Siden $ t_0, ..., t_n $ er forskjellige, så vil nevneren aldri være lik 0
    \item $ L_0, ..., L_n $ er polynomer av grad n
    \item $ L_k(t_{ \mathcal{l}}) = \delta_{k \mathcal{l}} for k, \mathcal{l} \in \left\{ 0, ..., n \right\}  $ 
\end{itemize}

\begin{propositionbox}{1.5.4}{}
La $ t_0, ..., t_n \in \mathbb{R} $ være $ n+1 $ forskjellige tall, og definer $ L_0, ..., L_n $. Da er $ \left( L_0, ..., L_n \right)  $ en basis for $ \mathcal{P}_n $ 
\end{propositionbox}

\begin{corollarybox}{1.5.5}{}
For $ n+1 $ forskjllige tall $ t_0, ..., t_n \in \mathbb{R} $ og tilhørende verdier $ y_0, ..., y_n \in \mathbb{R} $ finnes et unikt polynom $ p \in \mathcal{P}_n $ slik at $p(t_k) = y_k $ for $ k=0, ..., n $ 



\begin{proofbox}{1.5.5}{}
Vi bruker 1.5.3 for eksistens, og 1.5.4 for unikhet

Fra 1.5.3 for $ t_0, ..., t_n $ være $ n+1 $ forskjellige tall, så finnes det et polynom $ p_k :\mathbb{R} \rightarrow \mathbb{R} $ som er  slik at
\[
    p(t_k) = \prod_{j=0}^{n} \left( t_k - t_j \right) = \left( t_k - t_1 \right) ... \left( t_k - t_n \right) = y_k
\]
Dette danner en basis for $ \mathcal{P}_n $, og vi vet da at p eksisterer.

Uttrykk 
\[
    p(t_k) = \sum_{ j=0 }^{ n } L_j(t_k) y_k = \sum_{ j=0 }^{ n } \delta_{jk} y_k = y_k
\]
Uttrykk et nytt polynom $ q(t_k) = y_k $ for $ k=0, ..., n $ , kan vi fra 1.5.4 finne skalarer $ \alpha_{}^{}, ..., \alpha_{n}^{} $ slik at 
\[
    q = \alpha_{0}^{} L_0 + ...  + \alpha_{n}^{} L_n
\]. Da er

\[
    y_k = q(t_k) = \alpha_{k}^{} L_k = \alpha_{k}^{} \forall k=0, ..., n
\]
Altså er $ q = \sum_{ k=0 }^{ n } y_kL_k = p $ 

 
\end{proofbox}
\end{corollarybox}

\begin{corollarybox}{1.5.6}{}
La $ p \in  \mathcal{P}_n $ være et polynom med $ n+1 $ forskjellige røtter - det finnes $ n+1 $ forskjellige tall $ t_0, ..., t_n \in \mathbb{R}$ slik at $ p(t_0), ..., p(t_n)=0 $. Da er $ p=0 $.

\begin{proofbox}{1.5.6}{}
Nullpolynomet $ q(t) = 0 $ ligger i $ \mathcal{P}_n $ og oppfyller samtidig $ q(t_0) = ... = q(t_n) = 0 $. Av 1.5.5, så finnes det bare et polynom i $ \mathcal{P}_n $  medd denne egenskapen, så $ p=q $ 
\end{proofbox}
\end{corollarybox}

\subsection*{1.5.4 Legendre polynomer}
\addcontentsline{toc}{section}{Legendre polynomer}

\begin{align*}
    p_0(t) &= 1\\
    p_1(t) &= t\\
    p_k(t) &= \frac{ 2k -1 }{ k } t p_{k-1}(t) - \frac{ k-1 }{ k }p_{k-2}(t)
\end{align*}
for $ k=2, 3,.... $ for $ t \in \mathbb{R} $  

\begin{oppgavebox}{1.5.7}{}
Bis at de fire første legendre polynomene er gitt ved 
\begin{align*}
    p_0(t) &=1 \\
    p_1(t) &= t\\
    p_2(t) &= \frac{ 3 }{ 2 }t^2 - \frac{ 1 }{ 2 }\\
    p_3(t) &= \frac{ 5 }{ 3 }t^3 - \frac{ 3 }{ 2 } t
\end{align*}

\begin{proofbox}{1.5.7}{}
$ p_0, p_1 $ er opplagt
\begin{align*}
    p_2(t) &= \frac{ 2 \cdot \left( 2-1 \right)  }{ 2 } \cdot t \cdot \left( t \right) - \frac{ 1 }{ 2 } \cdot 1\\
    &= \frac{ 3 }{ 2 }t^2 - \frac{ 1 }{ 2 }
\end{align*}
$ p_3(t) $ er enkel å vise og
\end{proofbox}
\end{oppgavebox}

\begin{propositionbox}{1.5.8}{}
For enhver $ n \in \mathbb{N}_0 $ er $ p_n $ et n-te gradspolynom og $ \left( p_0, ..., p_n \right)  $ er en basis for $ \mathcal{P}_n $ 
\end{propositionbox}

\end{document}